\documentclass[12pt]{article}
\usepackage[T1]{fontenc}
\usepackage[utf8]{inputenc}
\usepackage{lmodern}
\usepackage{textcomp}
\usepackage{enumitem}
\usepackage{geometry}
\geometry{margin=1in}
\begin{document}
\begin{center}
Answer Key - Constitutional Law - Undergraduate Level Assessment 13 \
\small (Answers are ordered exactly as in the question paper.)
\end{center}

\section*{Multiple Choice}
\begin{enumerate}[label=\arabic*.]
\item \textbf{Answer:} D
\\ \textit{Options:} A) When the applicant is in the custody of the state; B) Where the state must investigate claims against it; C) The development of a legal framework to protect those within the state; D) To provide housing for all those homeless within a state
\\ \textit{Note:} Correct option: D - To provide housing for all those homeless within a state
\item \textbf{Answer:} B
\\ \textit{Options:} A) Fairness is a relative term.; B) The legal system is, in fact, unfair.; C) It sets a bad example.; D) The law is irrational and ambiguous.
\\ \textit{Note:} Correct option: B - The legal system is, in fact, unfair.
\item \textbf{Answer:} C
\\ \textit{Options:} A) There are fundamental differences between individual women.'; B) Men and women have different conceptions of the feminist project.'; C) Women look to context, whereas men appeal to neutral, abstract notions of justice.'; D) Men are unable to comprehend their differences from women.'
\\ \textit{Note:} Correct option: C - Women look to context, whereas men appeal to neutral, abstract notions of justice.'
\item \textbf{Answer:} D
\\ \textit{Options:} A) The language of the law is generally unclear.; B) Lawyers' arguments usually concern language.; C) Judges are prey to linguistic misunderstanding.; D) The rule of recognition cannot fully account for legal validity.
\\ \textit{Note:} Correct option: D - The rule of recognition cannot fully account for legal validity.
\item \textbf{Answer:} D
\\ \textit{Options:} A) Hobbes; B) Locke; C) Thomas Acquinas; D) Rousseau
\\ \textit{Note:} Correct option: D - Rousseau
\end{enumerate}

\section*{True / False}
\begin{enumerate}[label=\arabic*.]
\setcounter{enumi}{5}
\item \textbf{Answer:} False
\\ \textit{Options:} A) True; B) False
\\ \textit{Note:} LLM-generated false statement. Correct answer: 'Law and Economics'.
\item \textbf{Answer:} True
\\ \textit{Options:} A) True; B) False
\\ \textit{Note:} LLM-generated true statement based on MCQ.
\item \textbf{Answer:} True
\\ \textit{Options:} A) True; B) False
\\ \textit{Note:} LLM-generated true statement based on MCQ.
\item \textbf{Answer:} True
\\ \textit{Options:} A) True; B) False
\\ \textit{Note:} LLM-generated true statement based on MCQ.
\item \textbf{Answer:} True
\\ \textit{Options:} A) True; B) False
\\ \textit{Note:} LLM-generated true statement based on MCQ.
\end{enumerate}

\section*{Long Form Response}
\begin{enumerate}[label=\arabic*.]
\setcounter{enumi}{10}
\item \textbf{Answer:} Animals have inherent worth which entitles them to the absolute right to live their lives with respect and autonomy.. Explain why this is the correct answer and provide supporting evidence. Reference key facts or concepts that support this choice.
\\ \textit{Note:} Long-form derived from MCQ; award credit for stating the correct idea and giving supporting reasoning.
\item \textbf{Answer:} It rejects the Kantian concern with individual rights, equality, and justice.. Explain why this is the correct answer and provide supporting evidence. Reference key facts or concepts that support this choice.
\\ \textit{Note:} Long-form derived from MCQ; award credit for stating the correct idea and giving supporting reasoning.
\end{enumerate}

\vfill
\noindent\textit{End of Answer Key}
\end{document}